\documentclass[12pt]{article}
\usepackage[utf8]{inputenc}
\usepackage[T1]{fontenc}
\usepackage{hyperref}
\title{Panasuan Language Resources for K-12 Curriculum Development}
\author{Compiled by Claude Sonnet 5.5 (claude-sonnet-5-5) for Andr\'as Kornai}
\date{}
\begin{document}
\maketitle

\section{Introduction}
Panasuan (ISO 639-3: psn; Glottocode: pana1303) is an Austronesian language of the Seko branch of South Sulawesi, spoken on the border of West and South Sulawesi, Indonesia, mainly in the Kalumpang and Seko subdistricts. Estimates run from about 800 to 1,000 speakers. Searches found only general profiles and one audio item; no school materials were found.

\section{Primary Educational Resources}
No primers, textbooks or curricula were found.

\section{Linguistic and Reference Materials}
\begin{itemize}
\item Panasuan language (Wikipedia): \url{https://en.wikipedia.org/wiki/Panasuan_language}
\item Glottolog entry: \url{https://glottolog.org/resource/languoid/id/pana1303.bigmap.html}
\item Verbix endangered languages page: \url{https://docs.verbix.com/EndangeredLanguages/Panasuan}
\end{itemize}

\section{Religious and Community Resources}
\begin{itemize}
\item Joshua Project profile, Panasuan in Indonesia (translation reported as started): \url{https://joshuaproject.net/people_groups/14290}
\item Joshua Project language page for psn: \url{https://joshuaproject.net/index.php/languages/psn}
\end{itemize}

\section{Digital and Audio Resources}
\begin{itemize}
\item Global Recordings Network, Panasuan (audio Bible lessons reported): \url{https://globalrecordings.net/en/language/psn}
\end{itemize}

\section{Curriculum Development Context}
The language is used as a first language by its community but is not known to be taught. Any curriculum would need new materials, orthography decisions and community partnership; the audio lessons could be a starting point for recordings.

\end{document}